% Cycle de vie d'une entité : automate à 3 états.
\begin{tikzpicture}[
  ->, >=Stealth, thick,
  node distance = 4.5cm,
  every state/.style = {
    minimum size = 1.8cm,
    align = center,
    font = \small\ttfamily,
    draw = black,
    thick
  },
  alive/.style    = {state, fill=green!15},
  dying/.style    = {state, fill=orange!25},
  inactive/.style = {state, fill=gray!20}
]

\node[alive]    (alive)    {ALIVE};
\node[dying, right of=alive]    (dying)    {DYING};
\node[inactive, right of=dying] (inactive) {INACTIVE};

% Self-loop sur ALIVE (tick normal)
\path (alive) edge[loop above]
  node[font=\footnotesize, align=center]{tick normal\\(boids, manger, fuir)} (alive);

% ALIVE -> DYING : world_despawn
\path (alive) edge[bend left=15]
  node[above, font=\footnotesize, align=center]{world\_despawn() \\ (mangé / famine)} (dying);

% DYING -> INACTIVE : death_timer expiré
\path (dying) edge[bend left=15]
  node[above, font=\footnotesize, align=center]{death\_timer\\$\leq 0$} (inactive);

% INACTIVE -> ALIVE : respawn
\path (inactive) edge[bend left=45]
  node[below, font=\footnotesize, align=center]{respawn\_ready $\leq$ elapsed\\$\land$ alive\_count $<$ flock\_size} (alive);

\end{tikzpicture}
